% ---------------------------------------------------
\chapter{Limit Studies Discussion: Replication and Repair}
\label{chap:limit-discussion}
% ---------------------------------------------------

\section*{Overview}
This chapter consolidates the findings of all limit studies conducted for both
replication and repair mechanisms.  
The goal is to compare how reliability, overhead, and fault containment scale
across different fault-tolerance strategies and levels of granularity.
Replication mechanisms—Region-Level (RLCR), Function-Level (FLCR), and
Basic-Block-Level (BBLCR)—protect execution through redundancy and redirection.
Repair mechanisms, in contrast, reconstruct corrupted code or data in place,
relying on redundant encoding rather than multiple replicas.

Both methods operate under the same probabilistic framework introduced earlier:
survivability depends on the number of redundant elements, their validation
coverage, and the probability of an undetected fault.  
Together, replication and repair form the foundation of the proposed resilience
hierarchy—replication ensures continuity, while repair restores integrity.

% ---------------------------------------------------
\section*{Comparative Reliability}
The reliability data collected across all configurations reveal consistent
scaling behavior: as replication granularity increases, fault containment
improves and survivability extends to finer levels of program structure.
RLCR provides coarse fault isolation by selecting an entire verified image;
FLCR confines faults to the function level; and BBLCR achieves the highest
containment within individual control-flow blocks.

Repair mechanisms complement this trend by targeting individual instructions or
data words rather than entire structural units.  
Where replication avoids faults through redirection, repair mitigates them
through reconstruction.

\begin{figure}[t]
  \centering
  % \includegraphics[width=\linewidth]{Chapter-9/figs/reliability-trends.png}
  \caption{Placeholder for comparative reliability trends across replication and
           repair models (RLCR, FLCR, BBLCR, and Repair).}
  \label{fig:limit-reliability}
\end{figure}

% ---------------------------------------------------
\section*{Overhead and Structural Tradeoffs}
Replication and repair incur different types of cost.  
Replication introduces \textbf{static overhead} in the form of additional memory
for replicas and metadata structures (e.g., ITab, BTab), as well as
\textbf{dynamic overhead} from validation and redirection at runtime.
Repair mechanisms, meanwhile, add \textbf{computational overhead} for decoding
and reconstruction, especially when error density is high.

As granularity increases—from regions to basic blocks—the validation frequency
rises sharply, increasing both instruction count and runtime latency.
In contrast, repair overhead scales with the number of faults encountered rather
than program structure, making it more predictable in low-fault environments.

\begin{table}[h!]
  \centering
  \renewcommand{\arraystretch}{1.2}
  \begin{tabularx}{\linewidth}{|l|X|X|X|}
    \hline
    \textbf{Mechanism} & \textbf{Static Overhead} & \textbf{Dynamic Overhead} & \textbf{Notes} \\ \hline
    RLCR  & High (replica images) & Low (startup validation) & Coarse isolation \\ \hline
    FLCR  & Moderate (function replicas) & Moderate (call redirection) & Balance of cost and containment \\ \hline
    BBLCR & High (block metadata) & High (branch redirection) & Fine containment, high complexity \\ \hline
    Repair & Low (code parity) & Variable (reconstruction time) & Continuous protection \\ \hline
  \end{tabularx}
  \caption{Placeholder table comparing overhead characteristics for replication and repair mechanisms.}
  \label{tab:limit-overhead}
\end{table}

% ---------------------------------------------------
\section*{Fault Containment and Recovery Scope}
Each technique limits the spatial and temporal propagation of faults differently:
\begin{itemize}
  \item \textbf{RLCR} recovers at the region or image scale by selecting an alternate verified copy.
  \item \textbf{FLCR} confines corruption to individual functions, maintaining program flow through call-site indirection.
  \item \textbf{BBLCR} provides intra-function containment, validating and rerouting control at the level of basic blocks.
  \item \textbf{Repair} operates below these layers, reconstructing individual corrupted instructions or data words.
\end{itemize}

Replication therefore provides discrete containment boundaries, while repair
extends recoverability to the finest possible scale.
The two can operate in tandem—replication preserves execution continuity while
repair restores local integrity—forming a complementary defense structure.

\begin{figure}[t]
  \centering
  % \includegraphics[width=\linewidth]{Chapter-9/figs/containment-scope.png}
  \caption{Placeholder diagram illustrating fault containment scope for replication and repair mechanisms.}
  \label{fig:limit-containment}
\end{figure}

% ---------------------------------------------------
\section*{Combined Model: Repair within Replication}
In practice, replication and repair need not be exclusive.  
A hybrid model can use replication for immediate fault avoidance and repair for
deferred correction.  For example, if a replicated function fails validation, a
repair subsystem could reconstruct its corrupted copy and reintegrate it into
the validated pool.

This cooperative approach enables long-term reliability without unbounded
replication cost.  Replication provides runtime continuity, ensuring the system
always has a valid execution path, while repair reclaims lost replicas,
restoring redundancy over time.

\begin{figure}[t]
  \centering
  % \includegraphics[width=\linewidth]{Chapter-9/figs/hybrid-model.png}
  \caption{Placeholder for conceptual hybrid model showing interaction between replication and repair layers.}
  \label{fig:limit-hybrid}
\end{figure}

% ---------------------------------------------------
\section*{Summary}
The limit studies collectively demonstrate that increasing replication
granularity improves fault containment but at the expense of control complexity
and metadata size.  Repair mechanisms extend resilience beyond these limits by
addressing corruption directly, rather than redirecting around it.

Replication excels at maintaining continuous execution, while repair restores
full program integrity.  Together, they represent complementary dimensions of
software-based fault tolerance—replication for continuity, repair for recovery.

Future chapters will integrate these models under a unified runtime and compiler
framework to balance redundancy, cost, and self-healing capability.